\documentclass[11pt]{article}

\usepackage[margin=1in]{geometry}
\usepackage{amsmath, amssymb}

\newcommand{\todo}[1]{\textbf{[TODO: #1]}}
\newcommand{\coloneqq}{\mathrel{\vcentcolon=}}
\providecommand{\vcentcolon}{\mathrel{\mathop{:}}}
\newcommand{\xclean}{x_{\mathrm{clean}}}
\newcommand{\zclean}{z_{\mathrm{clean}}}
\newcommand{\aclean}{a_{\mathrm{clean}}}
\newcommand{\Dtheta}{D_\theta}
\newcommand{\Dref}{D_{\mathrm{ref}}}
\newcommand{\Etau}{\mathbb{E}_{\tau,\epsilon}}
\newcommand{\Dpref}{\mathcal{D}^{+}}
\newcommand{\Dneg}{\mathcal{D}^{-}}

\title{World-Model-Anchored Preference Optimization for Joint Flow-Matching Policies\\
\large \emph{Working draft}}
\author{Rooholla Khorrambakht \\ \emph{(co-authors TBD)}}
\date{\today}

\begin{document}
\maketitle

\begin{abstract}
We propose \emph{World-Model-Anchored Preference Optimization} (WAPO), a method for aligning a play-pretrained joint world-action transformer with small demonstration sets, without reward modeling and without distilling the policy into a tractable parametric family. Our backbone is a single causal flow-matching denoiser exposing independent per-frame noise levels for state and action, which lets the same network act as a world model, a policy, or a random-action proposer at inference time. We adapt the Diffusion-DPO bound to this dual-channel flow-matching setting, yielding a preference objective expressible directly through the model's denoising loss. Crucially, we use the frozen reference model itself --- in policy and random-action-proposer modes --- to synthesize the dispreferred trajectories needed for the DPO contrast, eliminating the dependence on external reward labels or human ranking. The aligned model is trained with low-rank adapters so the play-acquired prior is preserved by construction. We outline the method, derive its objective, describe its concrete realization in the \texttt{dreamerV4-UWM} codebase, and propose an experimental program built on the existing \texttt{pushT} suite plus a longer-horizon manipulation task. \todo{numbers, ablations}
\end{abstract}

\section{Introduction}
\label{sec:intro}

Joint world-action models trained on broad robotic ``play'' data offer an attractive prior for downstream control: they capture physical regularities and a wide span of behaviors, rather than collapsing onto a single task-specific policy. The dreamerV4-UWM line of work pushes this further by exposing \emph{independent} per-frame noise levels for the state and action streams of a flow-matching denoiser, so that one set of weights can be inference-time configured into three roles: world model, policy, or random-action proposer.

Two questions follow naturally from such a system:

\begin{enumerate}
    \item Once a play prior is in place, how do we shift it toward expert behavior on a target task using a small demonstration set, \emph{without} catastrophically forgetting the breadth that motivated play pretraining in the first place?
    \item What is the right post-training objective, given that the model's policy is implicit in a flow-matching denoiser whose action log-likelihood is intractable in closed form?
\end{enumerate}

A natural answer to (i) is preference optimization --- demonstrations are by construction preferred over rollouts of the broad prior --- and the natural answer to (ii) is the Diffusion-DPO formulation, which expresses the preference bound entirely in terms of the denoising loss. These two ingredients combine in a way that uses the model's distinguishing capabilities: \emph{the world model itself supplies the negative trajectories that the demonstrations are contrasted against.} The dual-channel architecture is no longer just an inference-time switch but a structural component of the alignment objective.

\paragraph{Contributions.}
\begin{itemize}
    \item A formulation of trajectory-level DPO for joint flow-matching world-action models with independent state/action noise levels, derived as a specialization of the Diffusion-DPO bound (\S\ref{sec:flow-dpo}).
    \item A self-contained source of dispreferred trajectories generated by the frozen reference model in its policy and random-action-proposer modes, removing the need for external reward labels (\S\ref{sec:wm-negatives}).
    \item A concrete instantiation in the \texttt{dreamerV4-UWM} codebase using LoRA adapters on the temporal blocks of the denoiser, with bit-equivalent fallback when alignment is disabled (\S\ref{sec:implementation}).
    \item A proposed experimental program that isolates each ingredient against ablation baselines, including a critical ``what does this do to modes 1, 2, and 3?'' regression check that we treat as a first-class outcome (\S\ref{sec:experiments}).
\end{itemize}

\section{Background}
\label{sec:background}

\subsection{Joint flow-matching world-action transformer}
\label{sec:bg-uwm}

The backbone (\texttt{dreamerv4uwm/models/dynamics.py:DreamerV4Denoiser}) is a causal transformer over a per-frame token sequence
\begin{equation}
    \mathrm{Frame}_t = [\, \mathbf{Z}_t,\ \mathbf{R}_t,\ \mathrm{IC}_t,\ \mathrm{AC}_t,\ \mathbf{A}_t \,],
\end{equation}
where $\mathbf{Z}_t$ are the latent tokens of frame $t$ produced by a frozen tokenizer, $\mathbf{R}_t$ is a small register, $\mathrm{IC}_t$ and $\mathrm{AC}_t$ are write-only \emph{control tokens} encoding the per-frame noise level for state and action respectively (and any shortcut step index), and $\mathbf{A}_t$ is the (possibly noisy) action token. Spatial attention is masked so the controls are read-only conditioning, and temporal attention is causal across frames.

\paragraph{Noise convention.}
We use the project's convention $\tau \in [0, 1]$ where $\tau = 1$ is clean data and $\tau = 0$ is pure noise, with linear forward mixing
\begin{equation}
    x_\tau \;=\; \tau\, \xclean \;+\; (1 - \tau)\,\epsilon, \qquad \epsilon \sim \mathcal{N}(0, I).
\end{equation}
The denoiser is trained as an $x$-predictor: given $x_\tau$ and $\tau$ it regresses $\xclean$. Independent per-frame noise levels $\tau_t^z$ for the state stream and $\tau_t^a$ for the action stream are encoded into the corresponding control tokens.

\paragraph{Three inference modes.}
A single trained denoiser is used in three roles by choosing the per-frame $(\tau^z_t, \tau^a_t)$ pattern:
\begin{itemize}
    \item \textbf{World model (WM):} clean context (states and actions), noisy horizon states, clean horizon actions $\Rightarrow$ predict horizon states given a candidate action sequence.
    \item \textbf{Policy ($\pi$):} clean context (states and actions), noisy horizon \emph{both} streams $\Rightarrow$ jointly generate states and actions on the horizon.
    \item \textbf{Random-action proposer (RAP):} \emph{noisy} context (states and actions), noisy horizon both streams $\Rightarrow$ broad joint generation that covers a search space for downstream planners.
\end{itemize}

\paragraph{Flow-matching loss.}
For a clean trajectory $\xi = (z_{1:T}, a_{1:T})$ and per-frame noise levels $\boldsymbol{\tau} = (\tau^z_{1:T}, \tau^a_{1:T})$, the per-sample denoising loss is
\begin{align}
    \ell_\theta(\xi, \boldsymbol{\tau}, \epsilon)
    \;=\; \sum_{t=1}^{T} \Big[\,
        & w(\tau^z_t)\,\big\| \Dtheta^z(z^{(\tau)}_{1:T}, a^{(\tau)}_{1:T}, \boldsymbol{\tau})_t - z_t \big\|^2 \nonumber \\
        + \;& w(\tau^a_t)\,\big\| \Dtheta^a(z^{(\tau)}_{1:T}, a^{(\tau)}_{1:T}, \boldsymbol{\tau})_t - a_t \big\|^2
    \,\Big],
    \label{eq:fm-loss-per-sample}
\end{align}
where $z^{(\tau)}, a^{(\tau)}$ are the streams mixed with independent Gaussian noise $\epsilon$ at the prescribed levels and $w(\cdot)$ is a noise-level loss weight (uniform or ramp; \S\ref{sec:weighting}). The expected loss is
$\mathcal{L}_{\mathrm{FM}}(\theta) = \mathbb{E}_{\xi, \boldsymbol{\tau}, \epsilon}\big[\ell_\theta(\xi, \boldsymbol{\tau}, \epsilon)\big].$

In the existing codebase, the joint $(\tau^z_{1:T}, \tau^a_{1:T})$ pattern is sampled per minibatch from a mixture indexed by mode --- \texttt{policy}, \texttt{wm}, \texttt{id}, \texttt{video}, \texttt{forcing} --- with mixing weights set by \texttt{train.mode\_weights}. We will reuse this sampler verbatim, but at \emph{evaluation} time within the alignment loop (\S\ref{sec:flow-dpo}).

\subsection{Direct Preference Optimization}
\label{sec:bg-dpo}

Given a reference policy $\pi_{\mathrm{ref}}$ and a dataset of preference pairs $(\xi^+, \xi^-) \sim \mathcal{D}$, DPO~\cite{rafailov2023direct} optimizes
\begin{equation}
    \mathcal{L}_{\mathrm{DPO}}(\theta)
    \;=\; -\,\mathbb{E}_{(\xi^+, \xi^-) \sim \mathcal{D}} \left[\,
        \log \sigma\!\left(\beta\,\big[\,r_\theta(\xi^+) - r_\theta(\xi^-)\,\big]\right)
    \right],
    \label{eq:dpo-base}
\end{equation}
where the implicit reward is $r_\theta(\xi) = \log \pi_\theta(\xi) - \log \pi_{\mathrm{ref}}(\xi)$ and $\beta > 0$ controls deviation from the reference. The objective is closed-form (no value network), but it presupposes that $\log \pi_\theta(\xi)$ is tractable.

\subsection{Diffusion-DPO}
\label{sec:bg-diffdpo}

For diffusion / flow-matching generative models, the trajectory log-likelihood is intractable but admits a variational upper bound expressible through the denoising loss. \emph{Diffusion-DPO}~\cite{wallace2024diffusion} substitutes this bound into Eq.~\eqref{eq:dpo-base}: up to a positive constant,
\begin{equation}
    \log \pi_\theta(\xi) - \log \pi_{\mathrm{ref}}(\xi)
    \;\geq\; -\,\Etau\!\left[\, w(\tau)\,\big(\|\Dtheta(\xi_\tau, \tau) - \xi\|^2 - \|\Dref(\xi_\tau, \tau) - \xi\|^2\big) \,\right]
    + C,
    \label{eq:diff-dpo-bound}
\end{equation}
which after substitution into Eq.~\eqref{eq:dpo-base} gives a fully tractable preference loss. We adapt this bound to the joint state-action setting in \S\ref{sec:flow-dpo}.

\section{Method}
\label{sec:method}

\subsection{Flow-DPO over joint state-action trajectories}
\label{sec:flow-dpo}

Define the per-trajectory \emph{advantage} of the aligned model over the reference under shared noise samples $(\boldsymbol{\tau}, \epsilon)$:
\begin{equation}
    \Delta_\theta(\xi; \boldsymbol{\tau}, \epsilon)
    \;\coloneqq\;
    \ell_{\mathrm{ref}}(\xi, \boldsymbol{\tau}, \epsilon) \;-\; \ell_\theta(\xi, \boldsymbol{\tau}, \epsilon),
    \label{eq:advantage}
\end{equation}
where $\ell_\theta$ is the per-sample flow-matching loss in Eq.~\eqref{eq:fm-loss-per-sample}. By Eq.~\eqref{eq:diff-dpo-bound}, $\mathbb{E}[\Delta_\theta(\xi;\cdot)]$ is a lower bound on $\log\pi_\theta(\xi)/\pi_{\mathrm{ref}}(\xi)$ up to a constant that cancels in differences.

The \textbf{Flow-DPO loss} for a preference pair $(\xi^+, \xi^-)$ is
\begin{equation}
    \boxed{\,
    \mathcal{L}_{\mathrm{FDPO}}(\theta)
    \;=\; -\,\mathbb{E}\!\left[\,
        \log \sigma\!\left(
            \beta\,\big[\Delta_\theta(\xi^+; \boldsymbol{\tau}^+, \epsilon^+) - \Delta_\theta(\xi^-; \boldsymbol{\tau}^-, \epsilon^-)\big]
        \right)
    \right],\,}
    \label{eq:flow-dpo}
\end{equation}
with the outer expectation taken over $(\xi^+, \xi^-) \sim \mathcal{D}$ and the noise variables. Two practical comments:

\paragraph{Shared noise across the pair.}
Following Wallace et al.~\cite{wallace2024diffusion}, we use the \emph{same} $(\boldsymbol{\tau}, \epsilon)$ for $\xi^+$ and $\xi^-$ within a pair. This is a strong variance-reduction trick: the constants in Eq.~\eqref{eq:diff-dpo-bound} match exactly between positive and negative, leaving only the model-dependent terms in the difference.

\paragraph{Mode sampler at evaluation time.}
We sample $(\boldsymbol{\tau}^z, \boldsymbol{\tau}^a)$ from the project's existing $\texttt{policy}$-mode schedule (clean context, sampled-$\tau$ horizon for both streams) when scoring a preference pair. This ensures the DPO gradient sees the policy regime --- the regime where $\xi^+$ vs.\ $\xi^-$ actually matters --- and does not leak weight onto the world-model regime (which we do not want to perturb). Other modes can be ablated.

\paragraph{Per-frame loss weighting.}
The loss weight $w(\tau)$ is the \texttt{ramp} or \texttt{uniform} schedule from \texttt{dreamerv4uwm/loss.py}. Note (\S\ref{sec:weighting}) that \texttt{ramp}, $w(\tau) = 0.9\tau + 0.1$, \emph{down-weights} the noisy regime under the project's $\tau{=}1{\Rightarrow}\mathrm{clean}$ convention. We default to \texttt{uniform} for FDPO since the preference signal at moderate-to-high noise carries the structural information about coarse trajectory shape.

\paragraph{Why this is honest about the architecture.}
The implicit policy in Eq.~\eqref{eq:flow-dpo} is the actual flow-matching denoiser $\Dtheta$: every gradient step adjusts the score field that the model uses to generate actions, not a parametric Gaussian projection. The dual-channel structure --- separate state and action terms in $\ell_\theta$ --- enters the loss linearly, so the alignment objective decomposes into ``align state predictions on $\xi^+$ vs.\ $\xi^-$'' plus ``align action predictions on $\xi^+$ vs.\ $\xi^-$''. We can mute either side via a coefficient if needed (\S\ref{sec:ablations}).

\subsection{Where do the negatives come from?}
\label{sec:wm-negatives}

The most fragile part of any trajectory-level DPO setup is the source of $\xi^-$. Random sequences make the contrast trivial; on-policy rollouts of $\pi_\theta$ require simulation; and human ranking is expensive. We exploit the dual-channel architecture to generate negatives from the \emph{frozen reference model itself}, in two complementary modes.

\paragraph{Reference policy negatives ($\xi^-_\pi$).}
Given an initial state $s_0$ matched to a demonstration, we run $\Dref$ in its \emph{policy} mode (clean context, jointly denoise horizon states and actions) using its standard sampling procedure. The resulting $\xi^-_\pi$ is what the broad play prior would produce from $s_0$. Demonstrations $\xi^+$ are by construction preferred to $\xi^-_\pi$: the demonstration is the expert sequence on this initial condition, while $\xi^-_\pi$ is a plausible play-prior continuation that may or may not be expert. The DPO contrast pushes $\Dtheta$ to assign higher implicit likelihood to demonstrations than to the prior on the same initial conditions.

\paragraph{RAP-mode negatives ($\xi^-_{\mathrm{RAP}}$).}
Generated by $\Dref$ in random-action-proposer mode (noisy context, joint generation), $\xi^-_{\mathrm{RAP}}$ samples are explicitly designed to be diverse and exploratory. They form a strong notion of ``physically plausible but task-uninformed'' trajectories. Including them in $\Dneg$ shapes $\Dtheta$ to assign implicit likelihood mass that is concentrated near demonstrations even relative to the prior's broad support.

\paragraph{Multi-source negative mixture.}
We use a mixture
\begin{equation}
    \Dneg \;=\; \alpha_\pi\,\Dneg_\pi \;\cup\; \alpha_{\mathrm{RAP}}\,\Dneg_{\mathrm{RAP}}
\end{equation}
with mixing weights $\alpha_\pi + \alpha_{\mathrm{RAP}} = 1$. The $\alpha$ ratio is treated as a hyperparameter (default $0.5/0.5$). Because both sources come from the \emph{frozen} $\Dref$, the negative distribution does not drift during training, eliminating one of the standard pathologies of online DPO.

\paragraph{Why the world-action backbone is load-bearing here.}
This is the key methodological point: the same network that we are trying to align is also the one that supplies the negatives. The dual-channel architecture is not decorative --- without RAP mode we lose the ``broad off-policy negatives'' arm, and without WM-quality joint generation we lose the ``in-distribution prior negatives'' arm. A purely policy-only baseline cannot replicate this without an external simulator.

\subsection{Segment-level preferences}
\label{sec:segments}

Trajectory-level Gaussian-style log-likelihoods are dominated by the longest streams; in our setting, the worst-positioned step of $\xi^+$ can flip the sign of the per-pair DPO term and inject high variance. We use \emph{segment-level} preferences: for each demonstration episode, we slide a window of length $L_s \leq T_{\mathrm{ctx}}$ across the episode and form preference pairs at each window position $\big(\xi^+_{[t,\,t+L_s]}, \xi^-_{[t,\,t+L_s]}\big)$ where the negative is generated by $\Dref$ from the demonstration's own context $[1,t-1]$. This:

\begin{enumerate}
    \item amplifies the number of supervisable preferences per demonstration roughly by $T_{\mathrm{episode}} / L_s$;
    \item avoids drift accumulation in the negative trajectory by always conditioning on the true demonstration history up to the segment start;
    \item matches the temporal scale at which the demonstration vs.\ prior comparison is well-defined --- whole-episode comparisons confound many local sub-decisions.
\end{enumerate}

We treat the segment length $L_s$ as a primary ablation knob (\S\ref{sec:ablations}).

\subsection{LoRA on the joint backbone}
\label{sec:lora}

The aligned model $\Dtheta$ is the reference $\Dref$ plus low-rank adapters on a chosen subset of layers. Concretely, for any linear projection $W \in \mathbb{R}^{d_{\mathrm{out}} \times d_{\mathrm{in}}}$ in the targeted set,
\begin{equation}
    W_\theta \;=\; W_{\mathrm{ref}} \;+\; \Delta W, \qquad \Delta W = B\,A, \quad A \in \mathbb{R}^{r \times d_{\mathrm{in}}},\ B \in \mathbb{R}^{d_{\mathrm{out}} \times r},\ r \ll \min(d_{\mathrm{in}}, d_{\mathrm{out}}),
\end{equation}
with $A$ Gaussian-initialized and $B$ initialized to zero so that $\Dtheta = \Dref$ at step zero. This makes the DPO term in Eq.~\eqref{eq:flow-dpo} exactly zero at initialization (the advantage difference vanishes), which gives a clean optimization start with no spurious early gradient.

We target the temporal blocks' $\{Q, K, V, O\}$ projections (the locus of multi-step decision-making) and leave the spatial blocks frozen by default. Spatial-LoRA and full-fine-tune are ablations.

\subsection{The full objective}
\label{sec:objective}

We optimize a convex combination of the DPO term and a small flow-matching anchor on the demonstration set:
\begin{equation}
    \mathcal{L}(\theta) \;=\; \mathcal{L}_{\mathrm{FDPO}}(\theta) \;+\; \lambda_{\mathrm{anchor}}\, \mathbb{E}_{\xi^+ \sim \Dpref}\!\left[\, \ell_\theta(\xi^+, \boldsymbol{\tau}, \epsilon) \,\right].
    \label{eq:full-loss}
\end{equation}
The anchor term is the same per-sample loss the backbone was trained with, restricted to demonstrations, and serves to keep the absolute scale of $\Dtheta$'s flow predictions calibrated to the data --- without it, DPO tends to inflate $\Dtheta$'s loss on $\xi^-$ by raising loss everywhere, which is a known failure mode~\cite{azar2024general,xu2024contrastive}. We default to small $\lambda_{\mathrm{anchor}} \in [10^{-2}, 10^{-1}]$ and ablate.

\section{Implementation in \texttt{dreamerV4-UWM}}
\label{sec:implementation}

This section makes the method concrete against the existing codebase. Code paths refer to current \texttt{main}.

\subsection{Stage 0: pretrained backbone}
We assume a checkpoint $\theta_{\mathrm{ref}}$ from \texttt{train\_dynamics\_uwm.py} on play data, ideally one of the recipes that already handles all three modes acceptably (e.g.\ the \texttt{progressive-forcing-action-video-wm-image} family in \texttt{docs/experiments.md}). The state of the backbone with respect to mode 3 is part of the experimental story --- a weaker mode-3 reference makes the RAP-mode negatives noisier but also makes the alignment-induced improvement easier to detect.

\subsection{Stage 1: precompute negatives}
\label{sec:impl-neg}
We add a one-time precomputation step that, for every demonstration window $\xi^+_{[t,t+L_s]}$, samples $K$ negatives from $\Dref$ in policy and RAP modes using the existing flow-matching sampler (\texttt{dreamerv4uwm/sampling.py}\todo{verify path}). Negatives are cached as a sibling shard layout (\texttt{shard\_XXXX\_neg.h5}) keyed on the same window indices, with payload \texttt{(image\_neg, action\_neg, source\_id, neg\_id)}. This keeps the alignment training loop offline-style, no on-policy sampling overhead.

This is computationally substantial but bounded: $|D^+|\cdot K\cdot 2$ rollouts per epoch-equivalent, with $K \in \{2, 4, 8\}$ in our setup. Each rollout is one forward pass of $\Dref$ at the chosen $T_{\mathrm{ctx}} + L_s$ horizon at modest sampling steps.

\subsection{Stage 2: alignment training script}

We add a new entry point \texttt{scripts/train\_align\_fdpo.py} parallel to \texttt{train\_dynamics\_uwm.py}. The script:
\begin{enumerate}
    \item Constructs $\Dref$ from a checkpoint, freezes it, places it in \texttt{eval} + \texttt{no\_grad}.
    \item Constructs $\Dtheta$ as a copy of $\Dref$ wrapped with LoRA adapters (PEFT-style, with $\Delta W = 0$ at init).
    \item For each step:
        \begin{enumerate}
            \item Draws a paired batch $(\xi^+, \xi^-)$ from a new \texttt{PreferenceShardedDataset} that joins the demonstration shards with the negative shards built in \S\ref{sec:impl-neg}.
            \item Samples one $(\boldsymbol{\tau}, \epsilon)$ pair via \texttt{UWMForwardProcess} restricted to the \texttt{policy} mode, applied identically to $\xi^+$ and $\xi^-$.
            \item Runs both $\Dtheta$ and $\Dref$ in a single batched forward (no gradient through $\Dref$).
            \item Computes $\ell_\theta(\xi^\pm, \boldsymbol{\tau}, \epsilon)$ and $\ell_{\mathrm{ref}}(\xi^\pm, \boldsymbol{\tau}, \epsilon)$ via a new \texttt{compute\_per\_sample\_uwm\_loss} that wraps \texttt{compute\_uwm\_loss} without reduction.
            \item Forms the FDPO term per Eq.~\eqref{eq:flow-dpo} and the anchor term per Eq.~\eqref{eq:full-loss}; sums and backpropagates only into the LoRA parameters.
        \end{enumerate}
    \item Periodically writes evaluation rollouts (mode 1, mode 2, mode 3) to verify no regression --- see \S\ref{sec:experiments}.
\end{enumerate}

\subsection{Code points to extend}
\label{sec:code-points}

\begin{itemize}
    \item \texttt{dreamerv4uwm/loss.py}: add \texttt{compute\_per\_sample\_uwm\_loss(...)} returning $(B,)$-shaped per-sample loss; refactor \texttt{compute\_uwm\_loss} to call into it. Existing call sites untouched.
    \item \texttt{dreamerv4uwm/models/dynamics.py}: add a thin \texttt{LoRAWrappedDenoiser} that takes a \texttt{DreamerV4Denoiser} and a target-module list and produces a wrapper exposing only the LoRA parameters as trainable. No structural change to the denoiser itself.
    \item \texttt{dreamerv4uwm/datasets.py}: \texttt{PreferenceShardedDataset} that pairs $(\xi^+, \xi^-)$ windows and shares the existing windowing math.
    \item \texttt{scripts/config/align/}: new config namespace mirroring \texttt{dynamics/} (e.g.\ \texttt{align/pushT-fdpo.yaml}) with $\beta$, $\lambda_{\mathrm{anchor}}$, $L_s$, LoRA rank, mode-mixture weights $\alpha_\pi/\alpha_{\mathrm{RAP}}$, and the source $\Dref$ checkpoint path.
    \item \texttt{hpc/slurms/align/pushT/*.slurm}: cloned from the dynamics slurms, swapping the entry point.
\end{itemize}

By construction, with $\beta = 0$ and $\lambda_{\mathrm{anchor}} = 1$ and LoRA rank $r$ but $\xi^- = \xi^+$, the script reduces to demonstration BC on the LoRA parameters --- providing a built-in sanity check.

\subsection{Compatibility with the shortcut variant}
\label{sec:shortcut-compat}

The shortcut bootstrap path (\texttt{train\_dynamics\_uwm\_with\_shortcut.py}) trains the same denoiser with an additional teacher-bootstrap head over a dyadic step grid. The FDPO loss is orthogonal: it operates on the flow branch ($\mathrm{step\_index} = \mathrm{max\_pow2}$) per Eq.~\eqref{eq:fm-loss-per-sample}. We therefore restrict the alignment loop to the flow branch and leave the bootstrap branch untouched, both at training and at inference. A future variant could attempt FDPO over the bootstrapped denoiser; we leave it out of scope.

\section{Proposed experiments}
\label{sec:experiments}

\subsection{Setup}
\label{sec:setup}

Pretrained reference $\Dref$: joint-model v1 (\texttt{checkpoints/dynamics/pushT/final/joint-model/v1/432360.pt}), trained on the pushT play dataset with the mode-mix recipe from \texttt{docs/experiments.md}. External reward model: Robometer, a language-conditioned per-frame progress predictor. The pushT prompt and Day-0 calibration audit are detailed in \S\ref{sec:reward-calibration}.

Window geometry: $T_{\mathrm{ctx}} = 1$ demo conditioning frame + $T_{\mathrm{hor}} = 9$ frames of sampled continuation per pair (10-frame windows). The short context is in-distribution for $\Dref$ because the pretraining recipe randomized the conditioning fraction.

\paragraph{Stage A pair generation (\S\ref{sec:wm-negatives} instantiated).}
For each demo window, $K=8$ rollouts are sampled in parallel from $\Dref$ via the progressive rolling sampler at $\texttt{context\_cond\_tau}=0.9$ (policy-mode siblings) or $\texttt{context\_cond\_tau}=0.5$ (RAP-mode siblings). Each rollout is decoded and scored by Robometer; the highest- and lowest-progress arms form $(\xi^+, \xi^-)$ if the spread exceeds $\texttt{MIN\_SPREAD}=0.10$. Smoke results on 50 windows: median spread $0.20$, drop rate $10\%$, per-arm progress span $\sim 0.4$ (\texttt{notebooks/explore\_fdpo\_pairs.ipynb}). Pos and neg now come from the same generative distribution and differ only in Robometer-scored progress --- this eliminates the source-confound (demo-vs-rollout) that drove v0/v0a DPO loss to saturation in $\sim 120$ steps.

\subsection{Claims}
\label{sec:claims}

We commit to three claims, each with a corresponding row in the headline table:
\begin{itemize}
    \item \textbf{C1 (source-confound fix).} Same-state sibling-pair preferences from $\Dref$ outperform demo-vs-rollout pairs (v0/v0a formulation) on the same DPO objective. The mechanism is that v0/v0a pairs are trivially separable by ``is this stream jittery'' (model rollouts) vs.\ ``is this stream smooth'' (human demos), which the model exploits as a free linear separator. Stage A pairs share a generative source and force the contrast onto task progress.
    \item \textbf{C2 (architecture-coupled negatives).} RAP-mode siblings (\texttt{context\_cond\_tau} $<1$) provide complementary contrast to policy-mode siblings, because the dual-channel architecture's RAP mode produces broader, more diverse rollouts that no single-channel policy can generate without an external mechanism. A mixture of the two sources outperforms either alone.
    \item \textbf{C3 (mode-retention).} Stage A LoRA fine-tuning improves Mode~2 (policy) task success while leaving Mode~1 (WM) loss and Mode~3 (RAP) entropy within $\pm 2\times$ $\Dref$. The unified backbone's other inference modes are preserved by construction --- a property that single-purpose diffusion-policy fine-tuning methods do not need to and do not satisfy.
\end{itemize}

Claims dropped from earlier drafts (for deadline scope): full-FT-vs-LoRA, segment-level-vs-episode, policy-only-DPO-vs-joint-DPO. Each was a reasonable knob but none load-bear the paper's thesis; each remains in scope for the follow-up.

\subsection{Headline table}
\label{sec:headline}

\begin{table}[h]
    \centering
    \small
    \begin{tabular}{lcccc}
        \hline
        \textbf{Method} & \textbf{Demo NLL $\downarrow$} & \textbf{Play NLL} & \textbf{Pair-acc.} & \textbf{Task success $\uparrow$} \\
        \hline
        $\Dref$ (no fine-tune)                  & --   & --   & --   & ?    \\
        BC-LoRA on demos                        & ?    & ?    & --   & ?    \\
        FDPO v0a (demo-vs-rollout)              & ?    & ?    & ?    & ?    \\
        \textbf{FDPO Stage A (policy sib., ours)} & \textbf{?} & \textbf{?} & \textbf{?} & \textbf{?} \\
        FDPO Stage A (RAP sib.)                 & ?    & ?    & ?    & ?    \\
        FDPO Stage A (mixed sib.)               & ?    & ?    & ?    & ?    \\
        \hline
    \end{tabular}
    \caption{Headline grid on pushT. Six runs, identical model, optimizer, $\beta$, anchor weight, LoRA rank, and step budget --- they differ only in the pair source. Row 4 vs.\ Row 3 isolates \textbf{C1} (Stage A vs.\ v0a source-confound baseline). Rows 5--6 vs.\ Row 4 isolate \textbf{C2} (architectural coupling). \emph{Pair-acc.} is the model's implicit preference margin on a held-out pair set (sanity that DPO is fitting the contrast at all).}
    \label{tab:headline}
\end{table}

Wall time: each FDPO row is $\sim 3$ h on local Blackwell at 5 epochs with the smoke-tested config (memory: 76.4 GB peak). BC-LoRA row is the same shape minus the DPO term. Six runs $\approx$ 18--24 h elapsed in tmux.

\subsection{Mode-retention table (first-class outcome)}
\label{sec:retention}

\begin{table}[h]
    \centering
    \small
    \begin{tabular}{lcccl}
        \hline
        \textbf{Method} & \textbf{Mode 1 (WM)} & \textbf{Mode 2 ($\pi$)} & \textbf{Mode 3 (RAP)} & \textbf{Verdict} \\
        & state-pred MSE & success rate & sample entropy & \\
        \hline
        $\Dref$ (frozen)                        & baseline & baseline & baseline & --   \\
        BC-LoRA                                 & ?        & ?        & ?        & ?    \\
        FDPO v0a                                & ?        & ?        & ?        & ?    \\
        FDPO Stage A (policy)                   & ?        & ?        & ?        & ?    \\
        FDPO Stage A (RAP)                      & ?        & ?        & ?        & ?    \\
        FDPO Stage A (mixed)                    & ?        & ?        & ?        & ?    \\
        \hline
    \end{tabular}
    \caption{Mode-retention regression check, reported alongside every row of Table~\ref{tab:headline}. A method that improves Mode~2 success but regresses Mode~1 MSE or Mode~3 entropy beyond $\pm 2\times$ $\Dref$ is reported as a wash (the verdict column says ``regression'' rather than ``win''). This is not supplementary; it is the property that distinguishes FDPO from single-purpose diffusion-policy fine-tuning (\S\ref{sec:related}).}
    \label{tab:retention}
\end{table}

v0 and v0a both regressed Mode 1/3 catastrophically (memory: play policy/wm $\sim$+1500 against ref) while improving Mode 2 nominally. Reporting Mode 2 alone would have hidden the failure mode and the v0/v0a learnings would have been lost. The retention table is therefore the gate, not a secondary metric.

\subsection{Reward-model calibration (Day-0 audit)}
\label{sec:reward-calibration}

Before Stage A pair generation, we verify that Robometer produces a meaningful ordered task-progress signal on pushT demo \emph{windows} (not full episodes; windows slice into arbitrary episode positions). The audit (\texttt{notebooks/explore\_fdpo\_pairs.ipynb}, cells 22--23) draws $N=16$ windows spaced across the dataset and checks three window-level conditions:
\begin{itemize}
    \item \emph{Within-window monotonicity}: fraction of windows with $\mathrm{progress}[-1] > \mathrm{progress}[0]$ exceeds $0.70$.
    \item \emph{Delta magnitude}: mean within-window increment $|\Delta|$ exceeds $0.05$, well above noise.
    \item \emph{Terminal-progress spread}: range across windows exceeds $0.30$ (Robometer resolves episode position).
\end{itemize}
PushT audit results: monotonicity $1.00$, mean $|\Delta| = 0.408$, terminal spread $0.559$. All three thresholds cleared by margin; Robometer is fit-for-purpose for sibling-rollout partitioning. The same audit will be applied to any second task.

\subsection{Secondary ablations}
\label{sec:ablations}

If time permits after the headline + retention tables fill, we report:
\begin{itemize}
    \item $\beta \in \{0.05, 0.1, 0.5\}$ around the chosen value, on the mixed-sibling row only.
    \item $\lambda_{\mathrm{anchor}} \in \{0.1, 0.5, 1.0\}$, same row.
    \item \texttt{MIN\_SPREAD} $\in \{0.10, 0.20\}$ pair-quality threshold.
    \item Per-stream FDPO mask: state-only (default per v1 finding) vs.\ both.
    \item Loss weighting $w(\tau)$: \texttt{uniform} (default) vs.\ \texttt{ramp}.
\end{itemize}
Deferred to follow-up: LoRA rank sweep, target-module sweep, $K$-arms scan, segment-length scan.

\subsection{Loss weighting note}
\label{sec:weighting}

The project's \texttt{ramp} weighting $w(\tau) = 0.9\tau + 0.1$ down-weights the noisy regime under the inverted $\tau{=}1{\Rightarrow}\mathrm{clean}$ convention. FDPO with \texttt{ramp} would put preference signal at $\tau \approx 1$ where the contrast is smallest. We default to \texttt{uniform} and ablate \texttt{ramp} once as a failure-mode characterization.

\section{Risks and mitigations}
\label{sec:risks}

\paragraph{Reward-hacking by raising absolute loss.} A standard DPO failure mode is to make $\xi^-$ much worse without making $\xi^+$ better, since only the difference enters $\sigma(\cdot)$. The anchor term in Eq.~\eqref{eq:full-loss} explicitly counteracts this by penalizing high absolute loss on demonstrations. We will monitor the absolute $\ell_\theta(\xi^+)$ over training; a sustained rise is a red flag.

\paragraph{Negative collapse.} If $\Dref$'s policy-mode samples are already close to $\xi^+$ for some initial conditions, the DPO contrast is near-zero and provides no gradient. We mitigate by (a) including RAP-mode negatives, which are deliberately broad, and (b) re-sampling negatives periodically (every $N$ epochs) to refresh the pool.

\paragraph{Distributional shift in the WM channel.} If FDPO drives the action-stream LoRA hard, the implicit conditioning on actions inside the WM mode could degrade rollout fidelity. Two safeguards: keep LoRA rank low; report Mode 1 RMSE as a hard regression gate (\S\ref{sec:regression-check}).

\paragraph{Compute cost of negatives.} Generating $K$ rollouts per window for the entire dataset is non-trivial but is a one-time precomputation per $\Dref$. For a refreshing schedule, we cap to one refresh per $\sim$10k optimization steps.

\paragraph{Episode boundary effects in segment preferences.} Sliding windows at the end of an episode hit the boundary; we follow the existing dataset windowing convention (\texttt{dreamerv4uwm/datasets.py}) and drop windows that span past the episode end.

\section{Related work (sketch)}
\label{sec:related}

\todo{Tighten with citations during paper writing.} Closest work: Diffusion-DPO~\cite{wallace2024diffusion}, which derives the bound we adapt; world-model policy optimization (WM-PO, Crossway diffuser) which uses world models to score actions but with explicit reward objectives; LoRA-based post-training~\cite{hu2022lora}; preference optimization for action diffusers in robotics. Distinct from prior robotics DPO in (a) joint state-action flow-matching backbone, (b) WM-self-generated negatives, (c) explicit retention guarantee on the world-model and random-action-proposer modes.

\section{Discussion}
\label{sec:discussion}

The aspect of this proposal we are most uncertain about is whether the \emph{magnitude} of the FDPO gradient through the temporal LoRA is large enough to meaningfully reshape the action distribution at sampling time, given that the implicit policy is realized through dozens of denoising steps. Diffusion-DPO works in image generation, where the analogous concern is real but empirically tractable; the robotics setting differs in that we care about the conditional density at very specific control-relevant frequencies, not perceptual statistics. The headline experiment to run before scaling up is therefore a sanity check that a small FDPO LoRA on a frozen $\Dref$ can recover the BC solution on a target task --- if it cannot, the bound is too loose at our scales and a different formulation is warranted.

The aspect we are most confident about is the negatives story: ``use the world-action model in its three inference modes to manufacture preference data'' is exactly the use case the dual-channel architecture was designed to enable, even if the original motivation was downstream planning rather than alignment.

\bibliographystyle{plain}
\begin{thebibliography}{9}

\bibitem{rafailov2023direct}
R.\ Rafailov, A.\ Sharma, E.\ Mitchell, S.\ Ermon, C.\ Manning, and C.\ Finn.
\newblock Direct Preference Optimization: Your Language Model is Secretly a Reward Model.
\newblock \emph{NeurIPS}, 2023.

\bibitem{wallace2024diffusion}
B.\ Wallace, M.\ Dang, R.\ Rafailov, L.\ Zhou, A.\ Lou, S.\ Purushwalkam, S.\ Ermon, C.\ Xiong, S.\ Joty, and N.\ Naik.
\newblock Diffusion Model Alignment Using Direct Preference Optimization.
\newblock \emph{CVPR}, 2024.

\bibitem{hu2022lora}
E.\ J.\ Hu, Y.\ Shen, P.\ Wallis, Z.\ Allen-Zhu, Y.\ Li, S.\ Wang, L.\ Wang, and W.\ Chen.
\newblock LoRA: Low-Rank Adaptation of Large Language Models.
\newblock \emph{ICLR}, 2022.

\bibitem{azar2024general}
M.\ G.\ Azar et al.
\newblock A General Theoretical Paradigm to Understand Learning from Human Preferences.
\newblock \emph{AISTATS}, 2024.

\bibitem{xu2024contrastive}
J.\ Xu et al.
\newblock Is DPO Superior to PPO for LLM Alignment? A Comprehensive Study.
\newblock \emph{arXiv:2404.10719}, 2024.

\todo{Add: Dreamer-V4; the dreamerV4-UWM tech report once written; flow-matching foundations (Lipman et al.); WM-PO; relevant robotics DPO work.}

\end{thebibliography}

\end{document}
